\section{Pricing Durability: A Configuration Ablation}
\label{sec:eval:dynamic:durability}

\autoref{sec:eval:dynamic:dirty} showed that the configuration used throughout
this chapter \emph{defers} the cost of persistence rather than avoiding it.
We disable the write-ahead log and periodic checkpointing so that \project{}
behaves as much like an in-memory system as its architecture permits.
This section measures what those mechanisms cost, by ablation: starting from
\project{} configured with both persistence mechanisms enabled, we remove one
at a time.

None of the intermediate configurations are used for any comparison reported in
this chapter; \textsc{none} is, and it is the fastest. The ablation eliminates
a hypothesis: if \project{}'s persistence machinery were expensive, the
throughput gap against Teseo, LiveGraph, and GraphOne (none of which are run
durably either) would need no further explanation. The measurements below show
that the machinery is cheap, which forces the explanation elsewhere.

\subsection{The configurations}
\label{sec:eval:dynamic:durability:configs}

\wt{} exposes two independent persistence mechanisms, a write-ahead log and
periodic checkpointing, which cross to give the four configurations in
\autoref{tab:durability-configs}, ordered by the strength of the guarantee each
provides.

The log alone provides a guarantee worth noting.
\project{} commits with \wt{}'s \texttt{sync=off}, which does not disable
syncing; it weakens the \emph{commit contract}. The redo record is buffered in
memory and the commit returns immediately, but the record is still forced to
disk when a log file fills and rotates, and again at every checkpoint. A crash
loses a \emph{bounded} suffix of recent commits rather than the whole run.
Describing this configuration as offering ``no durability'' is wrong; describing
it as offering the durability of a synchronous commit is equally wrong.

%%%%%%%%%%%%%%%%%%%%%%%%%%%%%%%%%%%%%%%%%%%%%%%%%%%%%%%%%%%%%%%%%%%%%%%%%%%%%%%%
\begin{table}[t]
  \centering
  \setlength{\tabcolsep}{4pt}\renewcommand{\arraystretch}{1.05}
  \small
  \begin{tabular}{l l l l}
  \toprule
  Configuration & Redo log & Checkpoints & A crash loses \\
  \midrule
  \textsc{wal-ckpt} & yes & yes & since the last rotation or checkpoint \\
  \textsc{wal}      & yes & no  & since the last log rotation \\
  \textsc{ckpt}     & no  & yes & since the last checkpoint \\
  \textsc{none}     & no  & no  & the entire run \\
  \bottomrule
  \end{tabular}
  \caption{The durability ladder, strongest guarantee first. Each row removes
  one guarantee from the row above it. \textsc{wal-ckpt} binds more tightly than
  \textsc{wal} despite an identical commit contract, because a checkpoint forces
  a synchronous flush of the log in addition to writing the data files.
  \textsc{none} is the configuration used elsewhere in this chapter.}
  \label{tab:durability-configs}
\end{table}
%%%%%%%%%%%%%%%%%%%%%%%%%%%%%%%%%%%%%%%%%%%%%%%%%%%%%%%%%%%%%%%%%%%%%%%%%%%%%%%%

\subsection{Method}
\label{sec:eval:dynamic:durability:method}

All cells insert the 260\,M edges of \texttt{graph500-24} into an initially
empty graph using the split edge-key layout with 72 writer threads, on the
machine described in \autoref{sec:eval:dynamic:setup}. The \wt{} cache is sized
at 800\,GiB so that the graph is resident throughout and eviction pressure never
binds; degree caching is disabled; memory-mapped I/O and statistics collection
are both off, as they perturb the ingest path. We interleave configurations
rather than running them in blocks, and report medians. Only the persistence
configuration varies between cells; we hold every other parameter, and the binary
itself, fixed.

\subsection{Results}
\label{sec:eval:dynamic:durability:results}

\autoref{tab:durability-ablation} gives the ablation. \project{}'s full
persistence machinery costs $1.67\times$ on ingest throughput: enabling both the
redo log and checkpointing takes the system from $1.70$ to $1.02$\,M edges/s.

%%%%%%%%%%%%%%%%%%%%%%%%%%%%%%%%%%%%%%%%%%%%%%%%%%%%%%%%%%%%%%%%%%%%%%%%%%%%%%%%
\begin{table}[t]
  \centering
  \setlength{\tabcolsep}{5pt}\renewcommand{\arraystretch}{1.05}
  \small
  \begin{tabular}{l l r r r}
  \toprule
  Configuration & Guarantee removed & Throughput & Gain over & Relative \\
                &                   & (M edges/s) & previous  & to \textsc{none} \\
  \midrule
  \textsc{wal-ckpt} & ---                         & 1.019 & ---       & $-40.1\%$ \\
  \textsc{wal}      & checkpoints                 & 1.071 & $+5.1\%$  & $-37.0\%$ \\
  \textsc{ckpt}     & redo log (checkpoints back) & 1.611 & $+50.5\%$ & $-5.2\%$  \\
  \textsc{none}     & checkpoints                 & 1.700 & $+5.5\%$  & --- \\
  \bottomrule
  \end{tabular}
  \caption{Durability ablation on \texttt{graph500-24}, 72 writers. Each row
  removes one guarantee from the row above; throughput rises monotonically as
  guarantees are given up.}
  \label{tab:durability-ablation}
\end{table}
%%%%%%%%%%%%%%%%%%%%%%%%%%%%%%%%%%%%%%%%%%%%%%%%%%%%%%%%%%%%%%%%%%%%%%%%%%%%%%%%

\paragraph{The redo log is the expensive mechanism; checkpointing is not.}
Almost the whole of the $1.67\times$ is attributable to the log. Removing it
(\textsc{wal}\,$\rightarrow$\,\textsc{ckpt}) is worth $50\%$, while the two
steps that add or remove checkpointing are worth about $5\%$ each, at the edge
of what these measurements resolve, and we do not claim the members of either
pair are distinguishable. This is a consequence of \project{}'s update
granularity. Each inserted edge is its own \wt{} transaction, so the log
receives one record per edge and its cost scales with the edge count, whereas a
checkpoint's cost scales with the volume of dirty pages it must reconcile and is
amortised over everything that accumulated since the last one.

\paragraph{Checkpointing is cheap at this scale, and only at this scale.}
The near-equality of \textsc{ckpt} and \textsc{none} should not be extrapolated.
\wt{}'s \texttt{checkpoint=(wait=$n$)} specifies the interval \emph{between} the
end of one checkpoint and the start of the next, not a period; as the graph
grows each checkpoint takes longer, so the duty cycle degrades over the life of
a run. In these ingests the observed interval stretches from roughly $78$\,s
early in the run to $479$\,s near its end. A \texttt{graph500-24} ingest
completes before that divergence has run its course, so the figures above
understate the steady-state cost of checkpointing on a longer or larger ingest.

\subsection{On synchronous commits}
\label{sec:eval:dynamic:durability:sync}

\wt{} can additionally be asked to make each commit durable before it returns,
rather than buffering the redo record. We measured this configuration and do not
report it as a rung of the ablation, because at \project{}'s transaction
granularity it does not measure what the ladder is measuring. \project{} commits
one edge per transaction, so a synchronous commit forces a durability wait per
edge ($260$\,M of them for the ingest above), and the resulting slowdown of
roughly $40\times$ is dominated by that granularity rather than by durability as
such. A system intending to ingest durably would batch many edges into each
transaction and amortise the wait across the batch; that is a different
experiment, and one we have not run. We note the measurement here so that its
absence from \autoref{tab:durability-ablation} is not mistaken for an oversight.

\subsection{Residual cost}
\label{sec:eval:dynamic:durability:residual}

\textsc{none} removes every persistence mechanism: no redo log, no
checkpointing, a cache large enough that the graph never approaches an eviction
threshold. Removing them was worth $1.67\times$ in total.
Yet \textsc{none} is the configuration behind every \project{} number reported
elsewhere in this chapter, including those where the gap against the in-memory
comparison systems is widest. The ablation therefore excludes the write-ahead
log and checkpointing as explanations for that gap.

The remaining cost cannot be switched off. Disabling the log and
the checkpoint server stops \wt{} from \emph{writing}, but leaves the
representation untouched: the graph is still held as page-oriented,
block-aligned B-tree images with byte-comparable keys, still reconciled when
pages are released, still governed by an extent allocator and an always-on
eviction protocol. \project{} pays for remaining \emph{persistable},
and \autoref{sec:eval:dynamic:dirty} shows the same thing from the other
direction: the deferred work does not disappear but accumulates. That cost is a property of the storage format, and no configuration
removes it.
